\documentclass{article}

\usepackage[preprint]{neurips_2024}
\usepackage[utf8]{inputenc}
\usepackage[T1]{fontenc}
\usepackage{microtype}
\usepackage{amsmath,amssymb,amsfonts}
\usepackage{booktabs}
\usepackage{graphicx}
\usepackage{subcaption}
\usepackage{multirow}
\usepackage[colorlinks=true,citecolor=blue,linkcolor=blue,urlcolor=blue]{hyperref}

\newcommand{\method}{Trajectory xg-SupCon}
\newcommand{\outputonly}{Output-only xg-SupCon}
\newcommand{\finalfeature}{Final-feature xg-SupCon}
\newcommand{\softre}{soft-rwERM}
\newcommand{\wg}{\textsc{WG-Acc}}
\newcommand{\avgacc}{\textsc{Avg-Acc}}
\newcommand{\nmi}{\textsc{NMI}}
\newcommand{\ari}{\textsc{ARI}}
\newcommand{\bacc}{\textsc{Bal-Acc}}
\newcommand{\minorityf}{\textsc{Minority-F1}}
\newcommand{\R}{\mathbb{R}}

\title{How Much Signal Is in Early Training Trajectories?\\A Matched-Budget Study of Pseudo-Group Inference}
\author{Anonymous Author(s)}
\date{}

\begin{document}
\maketitle

\begin{abstract}
Annotation-free robustness methods often depend more on pseudo-group inference quality than on the downstream robust objective itself. Motivated by recent work on early spurious-bias detection and training-dynamics analysis, we ask a narrower question: under matched warmup compute, do short early-training trajectories provide better pseudo-group signal than output-only early statistics or warmup-end features? We study a three-stage pipeline that (i) warms up a classifier for $T \in \{2,4\}$ epochs, (ii) infers within-class pseudo-groups from one of several signal families, and (iii) fine-tunes with either cross-group supervised contrastive learning or soft pseudo-group reweighted ERM. On Waterbirds, trajectory features slightly improve validation pseudo-group fidelity over the matched output-only control at $T=4$ (\nmi{} $0.321119 \pm 0.015119$ vs.\ $0.310790 \pm 0.019949$), but they reduce test worst-group accuracy under the same downstream learner ($0.491173 \pm 0.067613$ vs.\ $0.534787 \pm 0.055208$). Under the tighter $T=2$ budget, trajectories are worse on both fidelity and robustness, and the same ordering fails to transfer to soft pseudo-group reweighted ERM. Our main conclusion is therefore negative: within this matched-budget regime, modest fidelity gains from early trajectories do not reliably translate into better group robustness.
\end{abstract}

\section{Introduction}
Robustness to spurious correlations remains difficult when group annotations are unavailable. A common recipe is to infer pseudo-groups from model behavior and then train a downstream objective that emphasizes minority or cross-group structure. Recent papers argue that the critical bottleneck is often pseudo-group inference quality rather than the downstream optimizer alone \citep{han2024improving,chakraborty2024exmap,zheng2024selfguided}. At the same time, early training behavior appears to expose shortcut structure before the model fully converges \citep{yang2024spare,murali2023beyond}.

This paper asks a deliberately narrow question: \emph{at matched warmup cost, does richer early-training trajectory information yield better pseudo-group signal than output-only early statistics or warmup-end features?} The emphasis is on controlled comparison rather than on proposing a new robustness algorithm. We fix the downstream learner within each comparison, match warmup epochs across signal types, and isolate whether trajectory-specific descriptors help.

The answer in our experiments is mostly no. On Waterbirds, the trajectory signal improves validation \nmi{} only slightly relative to the matched output-only control at $T=4$ warmup ($0.321119$ vs.\ $0.310790$), yet test \wg{} drops from $0.534787$ to $0.491173$. At $T=2$, trajectory features are worse on both fidelity and robustness. A non-contrastive transfer test with soft pseudo-group reweighted ERM shows the same pattern: trajectory-based pseudo-groups have higher validation fidelity than output-only pseudo-groups (\nmi{} $0.321119$ vs.\ $0.310790$), but lower test \wg{} ($0.482866$ vs.\ $0.544652$). These results support a more cautious interpretation of early dynamics: some trajectory information is detectable, but in this budgeted regime it is not a reliable route to stronger group robustness.

Our contributions are:
\begin{itemize}
  \item We present a matched-budget benchmark for pseudo-group inference that compares warmup-end features, output-only early statistics, and richer multi-epoch trajectory descriptors under shared downstream learners.
  \item We introduce a lightweight trajectory-shape encoding and show that, on Waterbirds, it slightly improves validation pseudo-group fidelity at $T=4$ warmup but does not improve worst-group robustness; at $T=2$, it underperforms the output-only control on both metrics.
  \item We show that the ranking of pseudo-group signals does not transfer to a non-contrastive downstream learner on Waterbirds, and we provide a Colored MNIST sanity check, now including the ERM baseline, where richer trajectories improve clustering fidelity in an easier synthetic setting without surpassing the simple ERM reference on worst-group accuracy.
\end{itemize}

Section~\ref{sec:related} reviews prior work on pseudo-group inference and training dynamics. Section~\ref{sec:method} describes the matched-budget pipeline and the trajectory descriptors. Section~\ref{sec:experiments} presents the experimental setup, main results, ablations, and analysis. Section~\ref{sec:repro} summarizes reproducibility details, Sections~\ref{sec:discussion} and \ref{sec:conclusion} discuss limitations and conclude, and Appendix~\ref{sec:appendix} collects per-seed, stability, and scope-completeness details.

\section{Related Work}
\label{sec:related}
Our study sits at the intersection of annotation-free robustness, early spurious-bias detection, and training-dynamics analysis.

\textbf{Foundational spurious-correlation and group-robustness context.}
The modern annotation-free literature builds on a broader line of work on spurious correlations, environment shift, and worst-group risk. Group DRO established a strong supervised reference point by directly optimizing worst-group performance when group labels are available \citep{sagawa2020groupdro}. Human-annotation approaches showed that even relatively cheap side information can substantially improve robustness to spurious cues \citep{srivastava2020human}. In parallel, environment-invariance objectives such as IRM and sample-reweighting schemes such as Learning from Failure framed spurious correlations as a mismatch between easy shortcut features and more stable predictive structure \citep{arjovsky2019irm,nam2020lff}. We do not use group labels or environment annotations at training time, but this literature explains why pseudo-group quality is a central concern: once explicit groups are absent, some proxy for the worst-case structure has to replace them.

\textbf{Group inference and proxy-group construction.}
A broad line of work improves robustness to spurious correlations without group labels by inferring proxy groups, clustering within-class structure, or recovering latent concepts \citep{han2024improving,chakraborty2024exmap,arefin2024unsupervisedconcept,le2024out,you2024calibrating}. ExMap clusters explainability heatmaps from a trained model \citep{chakraborty2024exmap}, Unsupervised Concept Discovery clusters discovered latent concepts \citep{arefin2024unsupervisedconcept}, Out of Spuriousity extracts shortcut-related subnetworks \citep{le2024out}, and Calibrating Multi-modal Representations uses image-text alignment as an annotation-free grouping signal \citep{you2024calibrating}. \citet{han2024improving} make the most direct version of the pseudo-group-quality argument: better inferred groups are the key bottleneck for downstream robustness. We adopt that framing, but ask a narrower benchmarking question about which \emph{matched-cost warmup signal} is most useful for the inference stage.

\textbf{Objective-centric annotation-free robustness.}
Another branch of the literature focuses less on the proxy-group signal and more on the robust learner applied after or alongside it. Correct-N-Contrast uses contrastive learning to reduce shortcut reliance \citep{zhang2022correctncontrast}; Self-Guided Spurious Correlation Mitigation bootstraps proxy labels for robust training \citep{zheng2024selfguided}; and Evidential Alignment plus Elastic Representation propose annotation-free objectives that try to isolate class-relevant evidence \citep{ye2025improving,wen2025elastic}. These methods are important baselines for positioning because they show that better objectives can matter even when groups are unavailable. Our paper intentionally does not compete on objective novelty: we freeze the downstream learner as much as possible so that the comparison remains about pseudo-group signal quality.

\textbf{Early outputs, dynamics, and simplicity bias.}
\citet{yang2024spare} show that shortcut-related structure can be identified early in training and use early outputs to guide later importance sampling. That paper is the closest conceptual reference to our output-only control. A complementary diagnostic line studies training dynamics directly: \citet{murali2023beyond} show that spurious feature learning leaves characteristic temporal signatures, while \citet{qiu2024complexity} analyze how feature complexity influences what is learned in the presence of spurious correlations. Our study combines these perspectives by asking whether richer early dynamics are useful \emph{for group inference itself}, not only for analysis or reweighting.

\textbf{Contrastive representation learning.}
Supervised contrastive learning is now a standard way to improve class structure in representation space \citep{khosla2020supervised}, and more recent work explores disentangled or block-structured contrastive representations \citep{makino2025supervised}. We use xg-SupCon only as a controlled downstream consumer of pseudo-groups. The paper's claim is not that a new contrastive objective improves robustness, but that a matched-budget comparison can reveal when richer inferred structure is or is not actionable.

\textbf{Positioning.}
Unlike prior work, our contribution is neither a new standalone robustness algorithm nor a broad claim about subgroup discovery. It is a controlled empirical study of whether short early trajectories offer more useful pseudo-group signal than simpler matched-cost alternatives. The negative result is itself informative: higher pseudo-group fidelity does not necessarily imply higher worst-group accuracy.

\section{Method}
\label{sec:method}
\subsection{Problem setting and pipeline}
We consider binary classification with hidden spurious groups. Each sample has observed class label $y_i \in \{0,1\}$ and unobserved subgroup membership within class. The pipeline has three stages:
\begin{enumerate}
  \item Warm up a classifier for $T$ epochs, where $T \in \{2,4\}$.
  \item Infer within-class pseudo-groups from a chosen signal family.
  \item Fine-tune with a shared downstream learner.
\end{enumerate}

The central experimental variable is the pseudo-group signal. The downstream learner is fixed within each comparison so that changes in robustness can be attributed to pseudo-group quality rather than objective engineering.

\subsection{Warmup signal extraction}
For sample $i$ at warmup epoch $t$, let $z_i^{(t)}$ denote logits, $h_i^{(t)}$ the $\ell_2$-normalized penultimate feature, and $y_i$ the ground-truth class. Using two stochastic augmentations per sample at each checkpoint, we record
\begin{align}
  p_i^{(t)} &= \operatorname{softmax}(z_i^{(t)})_{y_i}, \\
  m_i^{(t)} &= z_i^{(t)}[y_i] - \max_{c \neq y_i} z_i^{(t)}[c], \\
  \ell_i^{(t)} &= \operatorname{CE}(z_i^{(t)}, y_i), \\
  u_i^{(t)} &= \operatorname{JS}(\operatorname{softmax}(z_{i,a}^{(t)}), \operatorname{softmax}(z_{i,b}^{(t)})), \\
  d_i^{(t)} &= 1 - \cos(h_i^{(t)}, h_i^{(t-1)}),
\end{align}
with $d_i^{(1)} = 0$. Here $p$ is true-class confidence, $m$ is margin, $\ell$ is per-sample cross-entropy, $u$ is augmentation disagreement, and $d$ is feature drift between adjacent epochs.

\subsection{Trajectory-shape encoding}
For any scalar sequence $x_i^{(1:T)}$, we compute
\begin{equation}
  g(x_i) = \big[x_i^{(1)}, x_i^{(T)}, \operatorname{mean}(x_i), \operatorname{std}(x_i), \Delta(x_i), \operatorname{slope}(x_i)\big],
\end{equation}
where $\Delta(x_i) = x_i^{(T)} - x_i^{(1)}$ and $\operatorname{slope}(x_i)$ is the least-squares slope over epoch index. We also define $t_i^{\mathrm{corr}}$ as the first epoch at which the sample is classified correctly with positive margin, or $T+1$ if this never occurs during warmup.

\subsection{Signal families}
All signals are standardized within class before clustering. If the within-class dimensionality exceeds $32$, we apply PCA to $16$ dimensions before fitting the mixture model.

\textbf{Trajectory signal.}
\begin{equation}
  s_i^{\mathrm{traj}} = [p_i^{(1:T)}, m_i^{(1:T)}, \ell_i^{(1:T)}, u_i^{(1:T)}, d_i^{(1:T)}, g(p_i), g(m_i), g(\ell_i), g(u_i), g(d_i), t_i^{\mathrm{corr}}].
\end{equation}

\textbf{Output-only control.}
\begin{equation}
  s_i^{\mathrm{out}} = [p_i^{(1:T)}, m_i^{(1:T)}, \ell_i^{(1:T)}, g(p_i), g(m_i), g(\ell_i), t_i^{\mathrm{corr}}].
\end{equation}
This matches the same warmup checkpoints and shape encoding but omits disagreement and feature-drift information.

\textbf{Final-feature control.}
\begin{equation}
  s_i^{\mathrm{feat}} = h_i^{(T)}.
\end{equation}

\subsection{Within-class pseudo-group inference}
For each class $y$, we cluster $\{s_i : y_i = y\}$ using a diagonal-covariance Gaussian mixture model with $K=2$ components. This choice matches the canonical majority/minority subgroup structure in our Waterbirds and Colored MNIST setups. Let $q_i \in \R^2$ denote the posterior mixture probabilities for sample $i$. We define a same-class cross-group disagreement score
\begin{equation}
  r_{ij} = 1 - q_i^\top q_j.
\end{equation}
The value is near zero when two examples are confidently assigned to the same pseudo-group and larger when their posterior assignments disagree.

\subsection{Downstream learners}
\textbf{Cross-group supervised contrastive learning.}
The main downstream learner combines cross-entropy with a weighted supervised contrastive term:
\begin{equation}
  \mathcal{L} = \mathcal{L}_{\mathrm{CE}} + \lambda \mathcal{L}_{\mathrm{xgSupCon}},
\end{equation}
where $\lambda = 0.5$. For an anchor $i$ with same-class positives $P(i)$,
\begin{equation}
  \mathcal{L}_{\mathrm{xgSupCon}}(i) =
  -\log
  \frac{\sum_{j \in P(i)} w_{ij} \exp(\operatorname{sim}(v_i, v_j) / \tau)}
       {\sum_{k \in B \setminus \{i\}} \exp(\operatorname{sim}(v_i, v_k) / \tau)},
\end{equation}
where $v_i$ is the normalized projection-head feature, $\tau = 0.07$, and
\begin{equation}
  w_{ij} = 0.1 + 0.9\, r_{ij}.
\end{equation}
We also evaluate a uniform-weight ablation that sets all same-class weights to one.

\textbf{Soft pseudo-group reweighted ERM.}
To test whether any signal ranking transfers beyond contrastive learning, we also fine-tune with weighted cross-entropy. For class $y$, let
\begin{equation}
  \pi_{y,k} = \frac{1}{|\{i : y_i = y\}|} \sum_{i: y_i = y} q_{i,k}
\end{equation}
be the estimated pseudo-group mass. Each example receives weight
\begin{equation}
  w_i^{\mathrm{ERM}} = \sum_k \frac{q_{i,k}}{\pi_{y_i,k} + 10^{-6}}.
\end{equation}

\section{Experiments}
\label{sec:experiments}
\subsection{Experimental setup}
\textbf{Datasets.}
Our main benchmark is Waterbirds, using the split summarized in \texttt{data\_summary.json}: $4795$ training examples, $1199$ validation examples, and $5794$ test examples. The training split is highly imbalanced, with a largest-to-smallest group ratio of $62.464286$, while validation and test have ratios $3.511278$ and $3.512461$. We also include Colored MNIST as a lightweight sanity check using the workspace implementation in \texttt{exp/colored\_mnist/}.

\textbf{Models and optimization.}
Waterbirds uses an ImageNet-pretrained ResNet-18 backbone; Colored MNIST uses a four-block convolutional network. Waterbirds experiments use batch size $64$, warmup for either $2$ or $4$ epochs, and downstream fine-tuning for $6$ epochs. Colored MNIST uses batch size $256$, $T=4$ warmup, and the same downstream epoch budget. The optimizer is AdamW with weight decay $10^{-4}$; the Waterbirds backbone learning rate is $3 \times 10^{-4}$ and all head learning rates are $10^{-3}$. For contrastive methods we use a two-layer projection head with output dimension $128$.

\textbf{Sampling, preprocessing, and hardware.}
Training uses balanced class batches. Waterbirds warmup applies random resized crops to $224 \times 224$, horizontal flips, color jitter, and random grayscale; evaluation resizes to $224 \times 224$. Colored MNIST uses padded random crops on $28 \times 28$ inputs and lighter color jitter. Experiments ran with Python $3.12.7$, PyTorch $2.10.0+\mathrm{cu}128$, torchvision $0.25.0+\mathrm{cu}128$, CUDA $12.8$, and a single NVIDIA RTX A6000 GPU, as recorded in \texttt{env\_report.json}.

\textbf{Metrics and variability.}
Our main downstream metrics are average accuracy, worst-group accuracy, and robustness gap. For pseudo-group fidelity, we report validation-set \nmi{}, \ari{}, minority-group F1 after the best binary cluster matching, and balanced accuracy of the matched pseudo-group assignments. We also report runtime and peak GPU memory. Waterbirds results use three seeds $\{7,17,27\}$; Colored MNIST uses one seed. Unless noted otherwise, every ``$\pm$'' entry in the Waterbirds tables and figures is the population standard deviation over those three seeds, matching the saved aggregate computation in \texttt{results\_summary.json}; Colored MNIST is reported as a single-seed sanity check with no uncertainty estimate.

\subsection{Main Waterbirds results}
Table~\ref{tab:main-waterbirds} summarizes the main Waterbirds comparison using values copied directly from \texttt{results\_summary.json}. Among the matched-budget xg-SupCon variants at $T=4$, the final-feature control reaches test \wg{} $0.535306 \pm 0.018269$, the output-only control reaches $0.534787 \pm 0.055208$, and the full trajectory variant reaches $0.491173 \pm 0.067613$. The trajectory signal therefore improves neither worst-group accuracy nor robustness gap, despite slightly higher validation fidelity. The same pattern appears across the extra fidelity metrics omitted in the earlier draft: relative to output-only at $T=4$, the trajectory signal improves \nmi{} from $0.310790 \pm 0.019949$ to $0.321119 \pm 0.015119$, \ari{} from $0.358425 \pm 0.021172$ to $0.403197 \pm 0.016493$, and \bacc{} from $0.799033 \pm 0.008468$ to $0.817585 \pm 0.006256$.

Two other patterns are notable. First, vanilla supervised contrastive learning without pseudo-groups is competitive in test \wg{} ($0.522326 \pm 0.025255$), which raises the bar for claims that richer pseudo-group signals are helping. Second, the soft reweighted ERM transfer test is also negative: output-only pseudo-groups give test \wg{} $0.544652 \pm 0.047062$, while trajectory pseudo-groups drop to $0.482866 \pm 0.056649$. Runtime and peak memory are nearly unchanged across pseudo-group variants, so the central difference is signal quality rather than resource usage.

\begin{table*}[t]
\caption{Waterbirds main results from \texttt{results\_summary.json}. Best value in each metric column is shown in \textbf{bold}. Higher is better for \avgacc{}, \wg{}, \nmi{}, and \minorityf{}, while lower is better for robustness gap, runtime, and peak GPU memory. Every entry is mean $\pm$ population standard deviation over seeds $\{7,17,27\}$.}
\label{tab:main-waterbirds}
\centering
\small
\resizebox{\textwidth}{!}{%
\begin{tabular}{lccccccc}
\toprule
Method & Test \avgacc{} $\uparrow$ & Test \wg{} $\uparrow$ & Test gap $\downarrow$ & Val \nmi{} $\uparrow$ & Val \minorityf{} $\uparrow$ & Runtime (min) $\downarrow$ & Peak memory (MB) $\downarrow$ \\
\midrule
ERM & 0.821022 $\pm$ 0.013077 & 0.518172 $\pm$ 0.052689 & 0.302849 $\pm$ 0.054106 & -- & -- & 3.502467 $\pm$ 0.011992 & \textbf{1045.165202 $\pm$ 2.029940} \\
Vanilla SupCon & 0.836613 $\pm$ 0.014410 & 0.522326 $\pm$ 0.025255 & 0.314287 $\pm$ 0.013238 & -- & -- & \textbf{3.482622 $\pm$ 0.008043} & 1046.517253 $\pm$ 0.117851 \\
\finalfeature{} ($T{=}4$) & 0.822633 $\pm$ 0.016705 & 0.535306 $\pm$ 0.018269 & 0.287326 $\pm$ 0.026512 & 0.314252 $\pm$ 0.015405 & 0.778111 $\pm$ 0.004139 & 3.635294 $\pm$ 0.000473 & 1046.600586 $\pm$ 0.000000 \\
\outputonly{} ($T{=}4$) & 0.843114 $\pm$ 0.010974 & 0.534787 $\pm$ 0.055208 & 0.308326 $\pm$ 0.045725 & 0.310790 $\pm$ 0.019949 & 0.791729 $\pm$ 0.005273 & 3.609370 $\pm$ 0.006701 & 1046.600586 $\pm$ 0.000000 \\
\method{} ($T{=}4$) & 0.846796 $\pm$ 0.012543 & 0.491173 $\pm$ 0.067613 & 0.355622 $\pm$ 0.055639 & \textbf{0.321119 $\pm$ 0.015119} & \textbf{0.816203 $\pm$ 0.004570} & 3.601641 $\pm$ 0.006173 & 1046.600586 $\pm$ 0.000000 \\
Output-only \softre{} ($T{=}4$) & 0.816362 $\pm$ 0.007614 & \textbf{0.544652 $\pm$ 0.047062} & \textbf{0.271710 $\pm$ 0.044166} & 0.310790 $\pm$ 0.019949 & 0.791729 $\pm$ 0.005273 & 3.547160 $\pm$ 0.006402 & 1046.600586 $\pm$ 0.000000 \\
Trajectory \softre{} ($T{=}4$) & \textbf{0.854217 $\pm$ 0.023802} & 0.482866 $\pm$ 0.056649 & 0.371351 $\pm$ 0.032996 & \textbf{0.321119 $\pm$ 0.015119} & \textbf{0.816203 $\pm$ 0.004570} & 3.554823 $\pm$ 0.006478 & 1046.600586 $\pm$ 0.000000 \\
\bottomrule
\end{tabular}
}
\end{table*}

\subsection{Matched-budget test}
The decisive comparison is the trajectory-versus-output-only test at matched warmup budget. Table~\ref{tab:matched-budget} shows that the trajectory signal is worse than the output-only control at both budgets on Waterbirds test \wg{}. At $T=2$, the trajectory variant trails by $0.067497$ in test \wg{} and by $0.011793$ in validation \nmi{}. At the same budget, it is only slightly better on the additional fidelity metrics (\ari{} $0.354703$ vs.\ $0.341944$; \bacc{} $0.797117$ vs.\ $0.792071$), which is not enough to improve robustness. At $T=4$, the trajectory variant gains $0.010330$ in \nmi{} and also improves \ari{} and \bacc{}, but still loses $0.043614$ in test \wg{}.

This matters because the original hypothesis predicted the opposite ordering: if trajectory information were especially useful under short compute, it should help most at $T=2$. Instead, the difference-in-differences goes the wrong way. Relative to output-only controls, the trajectory penalty in test \wg{} is larger at $T=2$ than at $T=4$ by $0.023884$.

\begin{table}[t]
\caption{Matched-budget Waterbirds comparison between output-only and trajectory signals. The final two columns report trajectory-minus-output differences at the same warmup budget.}
\label{tab:matched-budget}
\centering
\small
\resizebox{\linewidth}{!}{%
\begin{tabular}{lcccc}
\toprule
Method & Test \wg{} $\uparrow$ & Val \nmi{} $\uparrow$ & $\Delta$ \wg{} vs.\ output & $\Delta$ \nmi{} vs.\ output \\
\midrule
\outputonly{} ($T{=}2$) & 0.504673 $\pm$ 0.024364 & 0.295706 $\pm$ 0.020807 & 0.000000 & 0.000000 \\
\method{} ($T{=}2$) & 0.437175 $\pm$ 0.110919 & 0.283913 $\pm$ 0.023645 & -0.067497 & -0.011793 \\
\outputonly{} ($T{=}4$) & \textbf{0.534787 $\pm$ 0.055208} & 0.310790 $\pm$ 0.019949 & 0.000000 & 0.000000 \\
\method{} ($T{=}4$) & 0.491173 $\pm$ 0.067613 & \textbf{0.321119 $\pm$ 0.015119} & -0.043614 & 0.010330 \\
\bottomrule
\end{tabular}
}
\end{table}

\subsection{Ablations and analysis}
Table~\ref{tab:ablation} probes which components matter. Removing disagreement and drift recovers the output-only control. Removing the shape code lowers validation \nmi{} from $0.321119 \pm 0.015119$ to $0.308260 \pm 0.020994$ but raises test \wg{} from $0.491173 \pm 0.067613$ to $0.532191 \pm 0.068963$, again reinforcing that higher fidelity and better downstream robustness are decoupled in this regime. Uniform contrastive weights keep pseudo-group fidelity fixed but improve test \wg{} from $0.491173 \pm 0.067613$ to $0.510903 \pm 0.028066$, indicating that errors in how pseudo-groups are used downstream can dominate small fidelity differences.

Figures~\ref{fig:main-waterbirds} and \ref{fig:warmup} visualize the same story. Figure~\ref{fig:main-waterbirds} is restricted to the Waterbirds methods discussed in this section and reports three-seed error bars. Its left panel is a direct plot of the validation-\nmi{} values in \texttt{tables/main\_waterbirds\_table.csv} and \texttt{tables/matched\_budget\_waterbirds\_table.csv}; its right panel plots the corresponding test-\wg{} values from those same files. It shows a modest trajectory advantage in validation \nmi{} only at $T=4$, while output-only controls remain stronger on test \wg{} at both matched budgets. The warmup diagnostics figure illustrates why early trajectories may still appear promising: minority and majority examples can separate in confidence, loss, disagreement, and drift even after two warmup epochs. On Waterbirds, however, these signals are not yet reliable enough to improve the downstream robustness objective.

We also measured seed stability of pairwise disagreement scores on the Waterbirds validation set using \texttt{stability\_analysis.json}. Output-only pseudo-groups achieve mean pairwise Spearman correlation $0.519014 \pm 0.033269$ at $T=2$ and $0.553422 \pm 0.013578$ at $T=4$, while trajectory pseudo-groups achieve $0.514810 \pm 0.050027$ and $0.596938 \pm 0.027678$. The trajectory signal is therefore somewhat more stable at $T=4$, but that stability still does not translate into better \wg{}; Appendix~\ref{sec:appendix} shows the corresponding per-seed variability explicitly.

\begin{table}[t]
\caption{Waterbirds ablations and transfer checks at $T=4$. Every entry is mean $\pm$ population standard deviation over seeds $\{7,17,27\}$. The shape code improves validation pseudo-group fidelity, but none of the trajectory-based variants beats the output-only controls on test worst-group accuracy.}
\label{tab:ablation}
\centering
\small
\resizebox{\linewidth}{!}{%
\begin{tabular}{lccc}
\toprule
Variant & Val \nmi{} $\uparrow$ & Val \minorityf{} $\uparrow$ & Test \wg{} $\uparrow$ \\
\midrule
Full trajectory & \textbf{0.321119 $\pm$ 0.015119} & \textbf{0.816203 $\pm$ 0.004570} & 0.491173 $\pm$ 0.067613 \\
No disagreement/drift (output-only) & 0.310790 $\pm$ 0.019949 & 0.791729 $\pm$ 0.005273 & 0.534787 $\pm$ 0.055208 \\
No shape code & 0.308260 $\pm$ 0.020994 & 0.789429 $\pm$ 0.011411 & 0.532191 $\pm$ 0.068963 \\
Uniform SupCon weights & \textbf{0.321119 $\pm$ 0.015119} & \textbf{0.816203 $\pm$ 0.004570} & 0.510903 $\pm$ 0.028066 \\
Output-only \softre{} & 0.310790 $\pm$ 0.019949 & 0.791729 $\pm$ 0.005273 & \textbf{0.544652 $\pm$ 0.047062} \\
Trajectory \softre{} & \textbf{0.321119 $\pm$ 0.015119} & \textbf{0.816203 $\pm$ 0.004570} & 0.482866 $\pm$ 0.056649 \\
\bottomrule
\end{tabular}
}
\end{table}

\begin{figure*}[t]
\centering
\includegraphics[width=0.92\linewidth]{figures/waterbirds_main_comparison.pdf}
\caption{Waterbirds matched-budget summary with paper-level method names and three-seed error bars (population standard deviation), rendered from the aggregate rows in \texttt{tables/main\_waterbirds\_table.csv} and \texttt{tables/matched\_budget\_waterbirds\_table.csv}. Left: validation \nmi{}. Right: test \wg{}. The only clear fidelity gain for trajectories appears at $T=4$, while output-only controls remain stronger on worst-group accuracy at both warmup budgets. Final-feature clustering is competitive on \wg{} despite weaker fidelity than the full trajectory signal.}
\label{fig:main-waterbirds}
\end{figure*}

\begin{figure}[t]
\centering
\includegraphics[width=0.95\linewidth]{figures/waterbirds_warmup_diagnostics.pdf}
\caption{Warmup diagnostics pooled from \texttt{figures/warmup\_diagnostics\_data.json}, which was generated by replaying Waterbirds validation-split signal extraction for every saved output-only and trajectory run with warmup budgets $T \in \{2,4\}$ and seeds $\{7,17,27\}$. Confidence, margin, loss, disagreement, and drift show visible subgroup-dependent structure during warmup, which motivates trajectory-based pseudo-group inference. The main experimental result is that these signals are detectable but not sufficiently actionable to improve downstream worst-group accuracy under matched compute.}
\label{fig:warmup}
\end{figure}

\subsection{Colored MNIST sanity check}
Colored MNIST is easier and more synthetic than Waterbirds, so we use it only as a one-seed sanity check. Table~\ref{tab:colored-mnist} adds the previously omitted ERM reference. In that single run, the trajectory signal improves validation-set pseudo-group fidelity over the output-only control across all reported clustering metrics (\nmi{} $0.920212$ vs.\ $0.850755$, \ari{} $0.966703$ vs.\ $0.924016$, \bacc{} $0.976421$ vs.\ $0.945246$, minority-group F1 $0.985775$ vs.\ $0.964171$) and slightly improves test \wg{} ($0.838811$ vs.\ $0.834898$). However, the simple ERM baseline still achieves the best test \wg{} in this sanity check at $0.857590$. Because there is no multi-seed estimate of variability, we do not treat this as evidence comparable to the Waterbirds results. At most, it shows that richer trajectories can help clustering fidelity in a cleaner synthetic setting; it does not overturn the negative matched-budget result on Waterbirds.

\begin{table}[t]
\caption{Colored MNIST sanity check at $T=4$ using the single saved seed-$7$ run. ERM is included as a downstream-performance reference but has no clustering-based fidelity metrics because it does not infer pseudo-groups.}
\label{tab:colored-mnist}
\centering
\small
\resizebox{\linewidth}{!}{%
\begin{tabular}{lcccccc}
\toprule
Method & Test \avgacc{} $\uparrow$ & Test \wg{} $\uparrow$ & Val \nmi{} $\uparrow$ & Val \ari{} $\uparrow$ & Val \bacc{} $\uparrow$ & Val \minorityf{} $\uparrow$ \\
\midrule
ERM & \textbf{0.877400} & \textbf{0.857590} & -- & -- & -- & -- \\
\outputonly{} & 0.854800 & 0.834898 & 0.850755 & 0.924016 & 0.945246 & 0.964171 \\
\method{} & 0.840000 & 0.838811 & \textbf{0.920212} & \textbf{0.966703} & \textbf{0.976421} & \textbf{0.985775} \\
\bottomrule
\end{tabular}
}
\end{table}

\section{Reproducibility and Artifact Details}
\label{sec:repro}
The workspace includes the metadata needed to reproduce the paper. \texttt{env\_report.json} records the software and hardware environment; \texttt{data\_summary.json} records Waterbirds split sizes, class-group counts, imbalance ratios, and manifest hashes; \texttt{results.json} stores per-run outputs; and \texttt{results\_summary.json} stores the saved aggregates used for every numeric table cell in this paper.

Waterbirds uses seeds $\{7,17,27\}$ and Colored MNIST uses seed $7$. Each reported run directory under \texttt{exp/} contains \texttt{config.yaml}, \texttt{warmup\_log.csv}, \texttt{finetune\_log.csv}, \texttt{results.json}, and, for pseudo-group methods, \texttt{clustering\_outputs.json}. Warmup checkpoints are saved as \texttt{warmup\_epoch\_\{1,\dots,T\}.pt}. The primary reported Waterbirds artifacts live under \texttt{exp/waterbirds\_benchmark/}; the one-seed Colored MNIST artifacts live under \texttt{exp/colored\_mnist/}; and reported ablation runs live under \texttt{exp/waterbirds\_ablations/}. We explicitly exclude two smoke-test directories, \texttt{exp/\_smoke\_test\_waterbirds\_erm\_seed7/} and \texttt{exp/\_smoke\_test\_waterbirds\_out\_seed7/}, from the paper's artifact contract because they are one-epoch debugging runs rather than benchmark experiments.

Checkpoint selection follows the code in \texttt{exp/shared/core.py}. During warmup, the model state with the highest validation average accuracy is restored before signal extraction and downstream training. During downstream fine-tuning, the final reported validation and test metrics are computed from the checkpoint with the highest validation average accuracy over the six fine-tuning epochs. Pseudo-group fidelity is computed on the validation split only, using oracle validation groups for evaluation but never for training. Aggregate summaries in \texttt{results\_summary.json} use the arithmetic mean and population standard deviation across seeds, which is the uncertainty convention reported throughout this paper.

\section{Discussion and Limitations}
\label{sec:discussion}
The main lesson is that pseudo-group fidelity and downstream robustness can come apart. On Waterbirds, trajectory-based signals slightly improve validation clustering metrics at $T=4$, but they do not improve worst-group accuracy under either downstream learner. This weakens a simple ``better inferred groups imply better robustness'' narrative. In practice, small gains in pseudo-group separability may be too noisy, too class-conditional, or too poorly aligned with the downstream loss to matter.

Several limitations constrain how far we should generalize this result. First, the strongest evidence is on a single natural benchmark, Waterbirds, with a secondary one-seed Colored MNIST sanity check. Second, we only test short warmup budgets of $2$ and $4$ epochs and a fixed two-cluster assumption; the planned BIC-selected-$K$ sensitivity check was not run. Third, our pseudo-group fidelity metrics are computed with oracle validation groups and do not measure all downstream-relevant structure. Fourth, although the study is controlled, it does not exhaust alternative clustering models or downstream objectives. The appropriate claim is therefore narrow: within this particular matched-budget regime, modest gains in validation pseudo-group fidelity are not a reliable indicator of better worst-group accuracy. A stronger future study would broaden the dataset suite, vary the number of inferred groups, and investigate when fidelity metrics are predictive of robust accuracy.

\section{Conclusion}
\label{sec:conclusion}
We studied whether short early-training trajectories provide better pseudo-group signal than matched-cost output-only statistics or warmup-end features. The answer is mixed but mostly negative. On Waterbirds, the trajectory signal achieves slightly higher validation \nmi{} than the output-only control at $T=4$ ($0.321119$ vs.\ $0.310790$), yet it lowers test worst-group accuracy under xg-SupCon ($0.491173$ vs.\ $0.534787$) and under soft pseudo-group reweighted ERM ($0.482866$ vs.\ $0.544652$). At the tighter $T=2$ budget, trajectories are worse on both fidelity and robustness. Colored MNIST shows that richer trajectories can help in an easier setting, but the main matched-budget finding remains: modest improvements in inferred pseudo-group fidelity do not reliably translate into stronger group robustness.

More broadly, these results suggest that early dynamics should be treated as a diagnostic signal rather than assumed to be a reliable robustness intervention. For annotation-free group robustness, the next challenge is not only finding richer signals, but understanding which aspects of inferred group structure actually matter to downstream worst-group performance.

\begin{thebibliography}{99}

\bibitem[Arjovsky et~al.(2019)Arjovsky, Bottou, Gulrajani, and Lopez-Paz]{arjovsky2019irm}
Martin Arjovsky, L{\'e}on Bottou, Ishaan Gulrajani, and David Lopez-Paz.
\newblock Invariant risk minimization.
\newblock \emph{arXiv preprint arXiv:1907.02893}, 2019.

\bibitem[Arefin et~al.(2024)Arefin, Zhang, Baratin, Locatello, Rish, Liu, and Kawaguchi]{arefin2024unsupervisedconcept}
Md Rifat Arefin, Yan Zhang, Aristide Baratin, Francesco Locatello, Irina Rish, Dianbo Liu, and Kenji Kawaguchi.
\newblock Unsupervised concept discovery mitigates spurious correlations.
\newblock In \emph{Forty-first International Conference on Machine Learning}, 2024.

\bibitem[Chakraborty et~al.(2024)Chakraborty, Sletten, and Kampffmeyer]{chakraborty2024exmap}
Rwiddhi Chakraborty, Adrian Sletten, and Michael C. Kampffmeyer.
\newblock ExMap: Leveraging explainability heatmaps for unsupervised group robustness to spurious correlations.
\newblock In \emph{Proceedings of the IEEE/CVF Conference on Computer Vision and Pattern Recognition}, pages 12017--12026, 2024.

\bibitem[Han and Zou(2024)]{han2024improving}
Yujin Han and Difan Zou.
\newblock Improving group robustness on spurious correlation requires preciser group inference.
\newblock In \emph{Forty-first International Conference on Machine Learning}, 2024.

\bibitem[Khosla et~al.(2020)Khosla, Teterwak, Wang, Sarna, Tian, Isola, Maschinot, Liu, and Krishnan]{khosla2020supervised}
Prannay Khosla, Piotr Teterwak, Chen Wang, Aaron Sarna, Yonglong Tian, Phillip Isola, Aaron Maschinot, Ce Liu, and Dilip Krishnan.
\newblock Supervised contrastive learning.
\newblock In \emph{Advances in Neural Information Processing Systems 33}, pages 18661--18673, 2020.

\bibitem[Makino et~al.(2025)Makino, Park, Tagasovska, Kudo, Coelho, Huetter, Yao, Hoeckendorf, Leote, Ra, Richmond, Cho, Regev, and Lopez]{makino2025supervised}
Taro Makino, Ji Won Park, Natasa Tagasovska, Takamasa Kudo, Paula Coelho, Jan-Christian Huetter, Heming Yao, Burkhard Hoeckendorf, Ana Carolina Leote, Stephen Ra, David Richmond, Kyunghyun Cho, Aviv Regev, and Romain Lopez.
\newblock Supervised contrastive block disentanglement.
\newblock \emph{arXiv preprint arXiv:2502.07281}, 2025.

\bibitem[Nam et~al.(2020)Nam, Cha, Ahn, Lee, and Shin]{nam2020lff}
Junhyun Nam, Hyuntak Cha, Sungsoo Ahn, Jaeho Lee, and Jinwoo Shin.
\newblock Learning from failure: De-biasing classifier from biased classifier.
\newblock In \emph{Advances in Neural Information Processing Systems 33}, 2020.

\bibitem[Le et~al.(2024)Le, Schl{\"o}tterer, and Seifert]{le2024out}
Phuong Quynh Le, J{\"o}rg Schl{\"o}tterer, and Christin Seifert.
\newblock Out of spuriousity: Improving robustness to spurious correlations without group annotations.
\newblock \emph{arXiv preprint arXiv:2407.14974}, 2024.

\bibitem[Murali et~al.(2023)Murali, Puli, Yu, Ranganath, and Batmanghelich]{murali2023beyond}
Nihal Murali, Aahlad Puli, Ke Yu, Rajesh Ranganath, and Kayhan Batmanghelich.
\newblock Beyond distribution shift: Spurious features through the lens of training dynamics.
\newblock \emph{Transactions on Machine Learning Research}, 2023.

\bibitem[Qiu et~al.(2024)Qiu, Kuang, and Goel]{qiu2024complexity}
Guanwen Qiu, Da Kuang, and Surbhi Goel.
\newblock Complexity matters: Feature learning in the presence of spurious correlations.
\newblock In \emph{Proceedings of the 41st International Conference on Machine Learning}, pages 41658--41697, 2024.

\bibitem[Sagawa et~al.(2019)Sagawa, Koh, Hashimoto, and Liang]{sagawa2020groupdro}
Shiori Sagawa, Pang Wei Koh, Tatsunori B. Hashimoto, and Percy Liang.
\newblock Distributionally robust neural networks for group shifts: On the importance of regularization for worst-case generalization.
\newblock \emph{arXiv preprint arXiv:1911.08731}, 2019.

\bibitem[Srivastava et~al.(2020)Srivastava, Hashimoto, and Liang]{srivastava2020human}
Megha Srivastava, Tatsunori Hashimoto, and Percy Liang.
\newblock Robustness to spurious correlations via human annotations.
\newblock In \emph{Proceedings of the 37th International Conference on Machine Learning}, pages 9109--9119, 2020.

\bibitem[Wen et~al.(2025)Wen, Wang, Zhang, and Lei]{wen2025elastic}
Tao Wen, Zihan Wang, Quan Zhang, and Qi Lei.
\newblock Elastic representation: Mitigating spurious correlations for group robustness.
\newblock \emph{arXiv preprint arXiv:2502.09850}, 2025.

\bibitem[Yang et~al.(2024)Yang, Gan, Dziugaite, and Mirzasoleiman]{yang2024spare}
Yu Yang, Eric Gan, Gintare Karolina Dziugaite, and Baharan Mirzasoleiman.
\newblock Identifying spurious biases early in training through the lens of simplicity bias.
\newblock In \emph{Proceedings of The 27th International Conference on Artificial Intelligence and Statistics}, volume 238 of \emph{Proceedings of Machine Learning Research}, pages 2953--2961, 2024.

\bibitem[Ye et~al.(2025)Ye, Zheng, and Zhang]{ye2025improving}
Wenqian Ye, Guangtao Zheng, and Aidong Zhang.
\newblock Improving group robustness on spurious correlation via evidential alignment.
\newblock \emph{arXiv preprint arXiv:2506.11347}, 2025.

\bibitem[You et~al.(2024)You, Min, Dai, Sekhon, Staib, and Duncan]{you2024calibrating}
Chenyu You, Yifei Min, Weicheng Dai, Jasjeet S. Sekhon, Lawrence Staib, and James S. Duncan.
\newblock Calibrating multi-modal representations: A pursuit of group robustness without annotations.
\newblock In \emph{Proceedings of the IEEE/CVF Conference on Computer Vision and Pattern Recognition}, pages 26140--26150, 2024.

\bibitem[Zhang et~al.(2022)Zhang, Sohoni, Zhang, Finn, and R{\'e}]{zhang2022correctncontrast}
Michael Zhang, Nimit S. Sohoni, Hongyang R. Zhang, Chelsea Finn, and Christopher R{\'e}.
\newblock Correct-N-Contrast: A contrastive approach for improving robustness to spurious correlations.
\newblock \emph{arXiv preprint arXiv:2203.01517}, 2022.

\bibitem[Zheng et~al.(2024)Zheng, Ye, and Zhang]{zheng2024selfguided}
Guangtao Zheng, Wenqian Ye, and Aidong Zhang.
\newblock Learning robust classifiers with self-guided spurious correlation mitigation.
\newblock In \emph{Proceedings of the Thirty-Third International Joint Conference on Artificial Intelligence}, pages 5599--5607, 2024.

\end{thebibliography}

\appendix

\section{Additional Artifact-Backed Results}
\label{sec:appendix}

This appendix collects the additional details requested in review. Every number is copied from \texttt{results.json}, \texttt{results\_summary.json}, or \texttt{stability\_analysis.json}; no new experiments were run at paper-writing time.

\paragraph{Reported versus omitted experiment folders.}
The paper reports only three experiment families: benchmark runs in \texttt{exp/waterbirds\_benchmark/}, ablations in \texttt{exp/waterbirds\_ablations/}, and the single-seed sanity check in \texttt{exp/colored\_mnist/}. The folders \texttt{exp/\_smoke\_test\_waterbirds\_erm\_seed7/} and \texttt{exp/\_smoke\_test\_waterbirds\_out\_seed7/} are preliminary one-epoch validation runs used to verify the pipeline and are intentionally omitted from all tables and figures. Within \texttt{exp/waterbirds\_ablations/}, the saved notes \texttt{bic\_selected\_k\_optional/SKIPPED.md} and \texttt{no\_u\_d\_alias/SKIPPED.md} document, respectively, the skipped optional $K$-sensitivity study and the fact that the ``no $u/d$'' ablation is exactly the already-reported output-only control rather than a separate execution.

\begin{table*}[h]
\caption{Waterbirds validation pseudo-group fidelity metrics. Each entry is mean $\pm$ population standard deviation over seeds $\{7,17,27\}$ from \texttt{results\_summary.json}.}
\label{tab:appendix-fidelity}
\centering
\small
\resizebox{\textwidth}{!}{%
\begin{tabular}{lcccc}
\toprule
Method & Val \nmi{} $\uparrow$ & Val \ari{} $\uparrow$ & Val \bacc{} $\uparrow$ & Val \minorityf{} $\uparrow$ \\
\midrule
\finalfeature{} ($T{=}4$) & 0.314252 $\pm$ 0.015405 & 0.337524 $\pm$ 0.012970 & 0.790646 $\pm$ 0.006031 & 0.778111 $\pm$ 0.004139 \\
\outputonly{} ($T{=}2$) & 0.295706 $\pm$ 0.020807 & 0.341944 $\pm$ 0.024121 & 0.792071 $\pm$ 0.009915 & 0.791781 $\pm$ 0.000639 \\
\method{} ($T{=}2$) & 0.283913 $\pm$ 0.023645 & 0.354703 $\pm$ 0.036722 & 0.797117 $\pm$ 0.015104 & 0.801687 $\pm$ 0.005206 \\
\outputonly{} ($T{=}4$) & 0.310790 $\pm$ 0.019949 & 0.358425 $\pm$ 0.021172 & 0.799033 $\pm$ 0.008468 & 0.791729 $\pm$ 0.005273 \\
\method{} ($T{=}4$) & \textbf{0.321119 $\pm$ 0.015119} & \textbf{0.403197 $\pm$ 0.016493} & \textbf{0.817585 $\pm$ 0.006256} & \textbf{0.816203 $\pm$ 0.004570} \\
No shape code ($T{=}4$) & 0.308260 $\pm$ 0.020994 & 0.350102 $\pm$ 0.026152 & 0.796008 $\pm$ 0.011133 & 0.789429 $\pm$ 0.011411 \\
\bottomrule
\end{tabular}
}
\end{table*}

\begin{table*}[h]
\caption{Per-seed Waterbirds matched-budget results from \texttt{results.json}. These rows make the instability pattern at $T=2$ explicit: the trajectory condition attains the best single-seed result (seed $7$) and the worst two single-seed results (seeds $17$ and $27$).}
\label{tab:appendix-per-seed-matched}
\centering
\small
\resizebox{\textwidth}{!}{%
\begin{tabular}{lccccc}
\toprule
Method & Seed & Test \wg{} $\uparrow$ & Val \nmi{} $\uparrow$ & Val \ari{} $\uparrow$ & Val \bacc{} $\uparrow$ \\
\midrule
\outputonly{} ($T{=}2$) & 7 & 0.524922 & 0.266804 & 0.308049 & 0.778244 \\
\outputonly{} ($T{=}2$) & 17 & 0.518692 & 0.314937 & 0.362215 & 0.801001 \\
\outputonly{} ($T{=}2$) & 27 & 0.470405 & 0.305378 & 0.355569 & 0.796967 \\
\method{} ($T{=}2$) & 7 & 0.593458 & 0.250518 & 0.302790 & 0.775801 \\
\method{} ($T{=}2$) & 17 & 0.347352 & 0.302105 & 0.381876 & 0.808971 \\
\method{} ($T{=}2$) & 27 & 0.370717 & 0.299115 & 0.379444 & 0.806578 \\
\outputonly{} ($T{=}4$) & 7 & 0.459502 & 0.282859 & 0.328530 & 0.787073 \\
\outputonly{} ($T{=}4$) & 17 & 0.590343 & 0.328198 & 0.374821 & 0.805543 \\
\outputonly{} ($T{=}4$) & 27 & 0.554517 & 0.321313 & 0.371925 & 0.804483 \\
\method{} ($T{=}4$) & 7 & 0.425234 & 0.300949 & 0.380002 & 0.808769 \\
\method{} ($T{=}4$) & 17 & 0.584112 & 0.337347 & 0.416917 & 0.822640 \\
\method{} ($T{=}4$) & 27 & 0.464174 & 0.325062 & 0.412672 & 0.821346 \\
\bottomrule
\end{tabular}
}
\end{table*}

\begin{table*}[h]
\caption{Per-seed Waterbirds transfer and ablation results from \texttt{results.json}.}
\label{tab:appendix-per-seed-ablation}
\centering
\small
\resizebox{\textwidth}{!}{%
\begin{tabular}{lccccc}
\toprule
Method & Seed & Test \wg{} $\uparrow$ & Val \nmi{} $\uparrow$ & Val \ari{} $\uparrow$ & Val \bacc{} $\uparrow$ \\
\midrule
\finalfeature{} ($T{=}4$) & 7 & 0.518692 & 0.332906 & 0.353462 & 0.797792 \\
\finalfeature{} ($T{=}4$) & 17 & 0.560748 & 0.314670 & 0.337418 & 0.791104 \\
\finalfeature{} ($T{=}4$) & 27 & 0.526480 & 0.295179 & 0.321693 & 0.783041 \\
No shape code ($T{=}4$) & 7 & 0.507788 & 0.279275 & 0.314033 & 0.780639 \\
No shape code ($T{=}4$) & 17 & 0.626168 & 0.317184 & 0.361058 & 0.800736 \\
No shape code ($T{=}4$) & 27 & 0.462617 & 0.328322 & 0.375216 & 0.806650 \\
Uniform SupCon weights ($T{=}4$) & 7 & 0.507788 & 0.300949 & 0.380002 & 0.808769 \\
Uniform SupCon weights ($T{=}4$) & 17 & 0.478193 & 0.337347 & 0.416917 & 0.822640 \\
Uniform SupCon weights ($T{=}4$) & 27 & 0.546729 & 0.325062 & 0.412672 & 0.821346 \\
Output-only \softre{} ($T{=}4$) & 7 & 0.478193 & 0.282859 & 0.328530 & 0.787073 \\
Output-only \softre{} ($T{=}4$) & 17 & 0.580997 & 0.328198 & 0.374821 & 0.805543 \\
Output-only \softre{} ($T{=}4$) & 27 & 0.574766 & 0.321313 & 0.371925 & 0.804483 \\
Trajectory \softre{} ($T{=}4$) & 7 & 0.495327 & 0.300949 & 0.380002 & 0.808769 \\
Trajectory \softre{} ($T{=}4$) & 17 & 0.545171 & 0.337347 & 0.416917 & 0.822640 \\
Trajectory \softre{} ($T{=}4$) & 27 & 0.408100 & 0.325062 & 0.412672 & 0.821346 \\
\bottomrule
\end{tabular}
}
\end{table*}

\begin{table}[h]
\caption{Seed-stability analysis from \texttt{stability\_analysis.json}. Higher values indicate more consistent pairwise pseudo-group disagreement scores across seeds.}
\label{tab:appendix-stability}
\centering
\small
\begin{tabular}{lc}
\toprule
Method & Pairwise disagreement Spearman $\uparrow$ \\
\midrule
\outputonly{} ($T{=}2$) & 0.519014 $\pm$ 0.033269 \\
\method{} ($T{=}2$) & 0.514810 $\pm$ 0.050027 \\
\outputonly{} ($T{=}4$) & 0.553422 $\pm$ 0.013578 \\
\method{} ($T{=}4$) & \textbf{0.596938 $\pm$ 0.027678} \\
\bottomrule
\end{tabular}
\end{table}

\paragraph{Skipped \texorpdfstring{$K$}{K}-sensitivity check.}
The optional BIC-selected-$K$ ablation from \texttt{plan.json} was not run. The saved note in \texttt{exp/waterbirds\_ablations/bic\_selected\_k\_optional/SKIPPED.md} states that the required benchmark, transfer, and three-seed Waterbirds ablations exhausted the planned budget, and that the execution plan explicitly prioritized those required experiments before optional $K$-sensitivity work. We therefore do not make any claim that the conclusions are robust to alternative cluster counts beyond noting this unrun scope limitation.

\end{document}
